\documentclass[11pt]{article}
\usepackage[margin=27mm]{geometry}
\usepackage{amsmath,amssymb,amsthm,booktabs,hyperref}
\newtheorem{theorem}{Theorem}
\newtheorem{lemma}[theorem]{Lemma}
\newtheorem{proposition}[theorem]{Proposition}
\newtheorem{corollary}[theorem]{Corollary}
\newcommand{\Q}{\mathbb Q}\newcommand{\Z}{\mathbb Z}
\newcommand{\Tr}{\operatorname{Tr}}\newcommand{\Nm}{\operatorname{Nm}}
\title{Twelve contacts and a square-class normalization of the intrinsic pole}
\author{\(\Z\Q\) Diophantine Project: audit and continuation}
\date{19 November 2026}
\begin{document}\maketitle
\begin{abstract}
The ordinary-abc eleven-coordinate source can use its first height scale,
rather than its second or fourth power. A two-adic trace lemma supplies
the missing square class for integer witnesses. We also remove a known
rational square from the intrinsic-pole input before forming the finite
algebra. Together with automatic point and height guards and one obsolete guard,
these changes lower the degree bound to $6\,361\,985$. No pole parameter
or rationality witness is silently deleted; the count remains eleven.
The sufficient fixed shift is ineffective. All statements retain the
elliptic, trace-height, prime-pattern and intrinsic-pole premises of v63.
\end{abstract}

\paragraph{Source credit.}
The supplied constructions and notes audited here are by Gabe Conant with
Astra~6.2, in the research bundles
\emph{H10-quartic-constructions} and
\emph{H10-eleven-variable-constructions}.
The audit, corrections, and further improvements stated in this paper are
subsequent project work. The original source files and their provenance
remain separately preserved.

\section{Statement and source}

Use the project at commit \texttt{d3a8c22}, the ordinary-abc derivative
source v63, the full auxiliary charts v71, and even-valuation tests v75.
The positive factors of Gabe Conant with Astra~6.2 and positive balancing denominators are
as specified in \texttt{fewvars/POSITIVE.md} and \texttt{BALANCED.md}
in the supplied 18 November 2026 bundles. Ordinary rational abc is
assumed, together with precisely the normalized group, formal comparison,
trace-height, simultaneous prime-selection and intrinsic-pole inputs
of those papers. A finite circuit audit is not an audit of those inputs.

\begin{theorem}\label{thm:main}
For every sufficiently large fixed odd integer $B$, the finite
specification accompanying this paper gives a polynomial in eleven total
rational coordinates, including its input, with input projection $\Z$.
Its degree is at most $6\,361\,985$, its support is less than $2^{212}$,
and its integer content is one. The threshold for $B$ is ineffective.
The stored $B=3$ fixture has coefficient one-norm at most
$2^{25322112}$; it is not asserted to be a defining member.
\end{theorem}

\section{An exact square-class transfer}

For $e,v,s$ put
\[
 P_e(v,s)=e(e+v^2)^2s^4-e(e-v^2)^2(1-e)^2.
\]
\begin{lemma}\label{lem:transfer}
Suppose $e=dt^2$ with rational $t\ne0$. The substitution
\[
 v_{\rm old}=tv,\qquad Y_{\rm old}=t^3Y
\]
is a bijection between rational solutions of
\[
 Y^2=d(d+v^2)^2s^4-d(d-v^2)^2(1-e)^2
\]
and of $Y_{\rm old}^2=P_e(v_{\rm old},s)$.
It includes zero ordinates and preserves any density assertion under
the resulting invertible rational scaling of the parameter plane.
\end{lemma}
\begin{proof}
Substitution gives $P_{dt^2}(tv,s)$ equal to $t^6$ times the displayed
right side. Division by $t^6$ is legitimate. The inverse changes are
$v=v_{\rm old}/t$ and $Y=Y_{\rm old}/t^3$.
\end{proof}
This elementary identity is valid over every field of characteristic
different from two; the rational setting is needed for the later source.
It is a change of coordinates, not an additional realization assumption.

\section{The trace at two}

Write $K=\Q(\sqrt2)$, $\pi=\sqrt2$, and
\[
 D=11-6\sqrt2,\qquad E:y^2=x^3-Dx,\qquad
 P=\left(\frac{57-28\sqrt2}{8},\frac{-660+417\sqrt2}{32}\right).
\]
Normalize $v_\pi(\pi)=1$ in $K_\pi$, so $v_\pi(2)=2$.
\begin{lemma}\label{lem:trace}
For every nonzero integer $n$, put $T_n=\Tr_{K/\Q}x([n]P)$.
Then
\[
 n^2T_n-57/4\in16\Z_2,
 \qquad v_2(T_n)=-2-2v_2(n),\qquad T_n\in\Q_2^2.
\]
\end{lemma}
\begin{proof}
Let $z=-x/y$ be the formal parameter. Direct norms give
$v_\pi(x(P))=-6$, $v_\pi(y(P))=-9$, and $v_\pi(z(P))=3$.
The equation gives
\[
 x=z^{-2}+D/x=z^{-2}+O(z^2)
\]
with integral formal coefficients; equivalently
$x=(1+\sqrt{1+4Dz^4})/(2z^2)$, whose coefficients are integral
Catalan multiples of powers of $D$.
The integral formal group has $[m](z)=mz+O(z^5)$.
Indeed its coefficients are integral, and the automorphism
$(x,y)\mapsto(-x,iy)$ sends $z$ to $iz$; commutation with $[m]$
excludes every exponent not congruent to one modulo four. This argument
may be made after the integral extension adjoining $i$ and descends.
For odd $m$, substitution at $v_\pi(z(P))=3$ consequently gives
\[
 m^2x([m]P)-x(P)\in\pi^6\mathcal O_{K_\pi}.
\]

For a point $Q$ with $v_\pi(x(Q))<0$, the duplication formula gives
the exact identity
\[
 4x(2Q)-x(Q)
   =\frac{D(3x(Q)^2+D)}{x(Q)(x(Q)^2-D)}.
\]
As $D$ is a unit, its valuation is $-v_\pi(x(Q))>0$, and
$v_\pi(x(2Q))=v_\pi(x(Q))-4$. Apply this successively to
$Q=[2^jm]P$. Multiplying each difference by $4^jm^2$ and
telescoping proves
\[
 n^2x([n]P)-x(P)\in\pi^6\mathcal O_{K_\pi}
\]
for all positive $n$; negative $n$ have the same abscissa.
Now $\pi^6\mathcal O=8\mathcal O$, and
$\Tr\mathcal O=2\Z_2$. Taking traces proves the congruence.
It gives $4n^2T_n\equiv57\pmod {64}$, hence this last number is a
two-adic unit congruent to one modulo eight. The valuation and square
class follow.
\end{proof}

\begin{corollary}\label{cor:scale}
For integral $X=n\ne0$ and odd integral $B$, the number
\[
 M=(X^2+B)^4+T_n^4
\]
is a square in $\Q_2$.
\end{corollary}
\begin{proof}
The input $\beta=n^2+B$ is integral, whereas $v_2(T_n)\le-2$.
Thus $1+(\beta/T_n)^4\in1+2^8\Z_2\subset\Q_2^2$, and
$T_n^4$ is a square. This is a statement about the prepared integer
witnesses, not a claim that $M$ is a square for all rational base tuples.
\end{proof}

\section{The twelve-contact source}

Retain the square-abscissa root $A=[n]P$, $n\in\Z\setminus\{0\}$,
and the two full auxiliary point charts $H,U$. Set
\[
 \beta=X^2+B,\quad T=\Tr x(A),\quad M=\beta^4+T^4,\quad a=M,
 \quad x=a-z,\quad y=a+z.
\]
The twelve shifts are
\[
 \{-9,-8,-7,-5,-4,-3,-2,-1,6,10,11,12\}.
\]
They are distinct, nonzero, have sum zero, and no two are opposites.
Let $G_0(t)=\prod_c(t+c)$ and
\[
 F=G_0(x)G_0(y),\qquad
 H_a=\frac12\frac{\partial F}{\partial a},\qquad
 u=-\frac{12F}{H_a},\qquad w=\frac{u}{(1+u)^2}.
\]
Guard $F,H_a,u$ and all original point and test denominators.
The former guard $a^2-z^2\ne0$ is deleted: no formula divides by it,
and the stronger height estimate below does not use it.

The six raw tests are the four direct or conjugated reciprocal
abscissae and the cross and target errors of v63. For the first four
let $C$ be the raw reciprocal test; for the last two let $C$ be the
square of the raw error. Use
\[
 r_C=\frac{C^2}{u(1+u^2C^2)},\qquad
 W=\prod_C(1+r_C^2)=W_0+\sqrt2W_1,
 \quad V=\Tr(W^{-1}),\quad b=\frac{4w}{V^2},\quad e=7b.
\]
These are exactly the balanced positive factors of Gabe Conant with Astra~6.2. The first
four tests have even valuations at every good unramified place by
the equation of $E$; no global square assertion is needed for the
full auxiliary charts. Hence the same unique-minimum proof supplies
the old positive-odd extraction and negative-odd preparation.

Since $V=2W_0/\Nm(W)$,
\[
 e=7u\left(\frac{\Nm(W)}{(1+u)W_0}\right)^2.
\]
Apply Lemma~\ref{lem:transfer} with $d=7u$ and the square factor in
parentheses. The sole terminal equation is therefore
\[
 Y^2=7u(7u+v^2)^2s^4
       -7u(7u-v^2)^2(1-e)^2.                 \tag{1}
\]
The free coordinates are the input, six point coordinates, the slope,
the two pole parameters, and the encoder: eleven in total.

\section{Height soundness and an automatic denominator}

Use ordinary rational abc with $\epsilon=1/100$ in the intermediate
derivative-height estimate of v63. Its proof, before replacing the
larger slot height by half the product height, gives for any fixed
positive exponent $m$ and $a=M^m$
\[
 \log r^+(u)\ge(k-6-3\epsilon)\max(h(x),h(y))
                         -(2+\epsilon)h(M)-O(1).
\]
Here $r^+(u)$ is the positive odd radical outside a fixed exceptional
set containing the shift denominators, differences, nonzero pair sums,
and $2k$. Constants are uniform in $B$. The contact-pair noncollision needed by
the estimate holds here: $M\ge B^4$ and $B\ge3$ make every
$2M+c_i+c_j$ positive.
The proof also applies when $a^2-z^2=0$: then one slot is zero, its
contact values are nonzero constants, and the simple-prime argument
for the other slot is unchanged. The derivative denominator remains
guarded.

Since $x+y=2M$, the larger slot height is at least
$h(M)/2-O(1)$. With $k=12$ the result is
\[
 \log r^+(u)\ge\frac{39}{40}h(M)-O(1).       \tag{2}
\]
The original binary-power bound with $\kappa=1001/1000$ gives
\[
 \alpha E_h\le h(M)+O(1),\qquad
 \alpha=\frac{997}{1001},\qquad E_h=h(1:\beta:T).
\]
If $X\ne n$, the retained local extraction makes $r^+(u)$ divide
the reduced numerator of $X-n$, so
\[
 \log r^+(u)\le\tfrac12E_h+\log(1+|n|).
\]
The exact positive margin is
\[
 \frac{39}{40}\frac{997}{1001}-\frac12=\frac{18863}{40040}>0.
\]
The retained trace estimate $n^2\le A_0h(T)+J_0$ now bounds
$E_h$ uniformly in $B$. Since $E_h\ge h(\beta)\ge\log B$,
this is impossible for sufficiently large fixed $B$. Thus $X=n$.

There is a second use of (2) which does not use the pole equation or
the desired conclusion $X=n$. For every rational tuple with $F,H_a$
nonzero, $M\ge B^4$. Thus for sufficiently large $B$, the right
side of (2) is positive. Consequently $u$ has a positive odd
valuation and cannot equal $-1$. The denominator $1+u$ is therefore
automatic. This is an eventual assertion under abc, not a claim about
every tuple in the $B=3$ fixture. Also $M>0$ at every rational tuple.
These facts justify the two additional norm-guard deletions in the
finite-algebra specification without circular use of the pole.

\section{Complete integer witnesses}

Fix $X=n\ne0$ and $A=[n]P$. Take the original multiplication family
$U=[n]H$ in the full auxiliary domain. The new scale $a=M$ is fixed
and positive before choosing the slope. It exceeds all shift
magnitudes for $B\ge3$, so the 24 rational positions $\pm(a+c)$
are distinct. All other exceptional sets, point neighborhoods and
reduction-order constraints are fixed in the old order.

At fixed finite and real infinity in the slope, the derivative quotient
satisfies $u\sim z^2/a$ and $w\sim a/z^2$. The fixed nonzero tests
stay bounded away from zero in the prepared neighborhoods, so the
balanced detector satisfies $W\to1$ and $V\to2$.
Corollary~\ref{cor:scale} makes $a$ a square over $\Q_2$, which is
precisely the square-class requirement formerly supplied by $a=M^2$
or $M^4$. Arbitrarily large negative $v_2(z)$ gives the required
positive even depth. At seven and the other fixed places, the same
large-slope preparation gives the required nonnegative depths.

Apply the retained simultaneous prime-pattern theorem to the 24
fixed positions. At every selected prime one contact has valuation
one, so $u$ has valuation one. The source compatible index refinement
makes all successful elliptic tests sufficiently deep; every positive
factor has residue one there, while the opposite inverse factors
are zero or one. Thus the reciprocal trace is a unit and the actual
pole input has positive odd valuation at the selected primes.
At every other prime the source denominator argument is unchanged:
negative odd valuations of $u$ give $W\equiv1$ at both slots, while
even valuations stay even after division by the rational square trace.
The finite support of the now-fixed scale is prepared in advance.
The actual input has square reduced denominator and the required
real interval $0<e<1$ and intrinsic local square class.

The intrinsic-pole converse supplies the old parameters and ordinate.
The square factor in Lemma~\ref{lem:transfer} is nonzero, so they give
rational parameters for (1). The density/avoidance conclusion transfers
by that invertible scaling. Every whole-character condition retained
by the compiler is handled by the same prepared finite avoidance;
the positive balancing and detector factors stay nonzero on rational
characters. The final factor $X$ includes the integer zero exactly
as in v63. This proves completeness and Theorem~\ref{thm:main}.

\section{The finite specification and audit}

The three point blocks retain rank $4^3=64$ and the terminal square
has rank two, for total rank 128. The point-prefix separator and
zero-terminal-root theorem are the audited constructions of v40 and
the notes of Gabe Conant with Astra~6.2. Four point norm atoms may also be omitted.
For the root chart $x=g^2$, $j^2=g^4-D$, the norm-root square
$\Nm(g^4-D)$ never vanishes at a rational base: vanishing would force
$(\Nm g)^4=\Nm D=49$, impossible in $\Q$.
For each full auxiliary chart, the retained ordinate norm forces
$\Nm(x^3-Dx)\ne0$ and hence $\Nm x\ne0$.
These facts restore point etaleness before character extraction.
After a rational character is obtained, $g=0$ would give $j^2=-D$,
which total positivity of $D$ excludes. The integral algebra and
separator are defined before these rational denominators are evaluated.
The four deleted atoms save total guard degree 14.

Guard recovery then proceeds through
the retained point/test denominators; positivity of $M$ and the
independent height proof that $1+u\ne0$; positivity of the six balanced
denominators; positivity of $W_0$ and $\Nm(W)$; and the invertible
intrinsic-pole scaling. No deleted guard is used to justify itself.

The final polynomial is
\[
 X\det(I-YG_*\,\mathrm{Mult}(\Omega)).
\]
The final guard degree is 650 and the separator operator degree is
49052, so
\[
 \deg\le1+128(1+650+49052)=6\,361\,985.
\]
Setting $Y=0$ gives $X$, proving integer content one. The support
bound counts total degree at most $d=6361984$ and encoder degree at
most 128 before multiplication by $X$:
\[
 \binom{d+11}{11}-\binom{d-128+10}{11}<2^{212}.
\]
The numerical coefficient bound concerns only the stored fixture.

The independent checker reconstructs all twelve factors and their
derivative, the scale and raw tests, all six balanced positive factors,
all 164 denominator records, and the new terminal equation. It compares
the point charts and original tests against the pinned v75 source and
replays all square extractions, guard degrees, envelopes and coefficient
majorants. It separately verifies the intrinsic square-factor identity,
the exact duplication identity, the fixed-point valuations and the
strict height margin. The proof of Lemma~\ref{lem:trace} covers every
integer $n$; it is not inferred from a finite list of multiples.

\begin{thebibliography}{9}\raggedright
\bibitem{source} \(\Z\Q\) Diophantine Project, v63 derivative detector,
v71 full auxiliary chart, v75 even-valuation detector; repository
\url{https://github.com/JoseBrox/diophantine-z-q_improvements},
commit \texttt{d3a8c22}. The relevant full sources accompany this bundle.
\bibitem{algebra} Project papers on unit norms v28, subalgebra norms v30,
and possibly vanishing terminal quadratic roots v40.
\bibitem{coauthor} Gabe Conant with Astra~6.2, research bundles of 18 November 2026,
\texttt{POSITIVE.md}, \texttt{BALANCED.md}, and
\texttt{counts/poles/ordinary\_scale/proof.md}.
\end{thebibliography}
\end{document}
